% Standalone problem document, generated from the adjacent README.md.
% Compile with XeLaTeX. Mathematical content is maintained in Markdown.
\documentclass[11pt,a4paper]{article}
\usepackage[fontsize=10.5pt]{scrextend}
\usepackage[margin=22mm,headheight=15pt,headsep=9mm,footskip=11mm]{geometry}
\usepackage{fontspec}
\setmainfont{texgyrepagella}[Extension=.otf,UprightFont=*-regular,BoldFont=*-bold,ItalicFont=*-italic,BoldItalicFont=*-bolditalic]
\setsansfont{texgyreheros}[Extension=.otf,UprightFont=*-regular,BoldFont=*-bold,ItalicFont=*-italic,BoldItalicFont=*-bolditalic]
\setmonofont{texgyrecursor}[Extension=.otf,UprightFont=*-regular,BoldFont=*-bold,ItalicFont=*-italic,BoldItalicFont=*-bolditalic,Scale=MatchLowercase]
\usepackage{amsmath,amssymb}
\usepackage{unicode-math}
\setmathfont{texgyrepagella-math.otf}
\usepackage{xcolor,microtype,enumitem,needspace,fancyhdr}
\usepackage{bookmark}
\definecolor{ink}{HTML}{17344B}
\definecolor{accent}{HTML}{197D80}
\definecolor{muted}{HTML}{596773}
\hypersetup{colorlinks=true,linkcolor=accent,urlcolor=accent,pdftitle={IE-20 - Precision
required for conjugate gradients to attain backward accuracy in n
steps},pdfauthor={Open Problems in Numerical Linear Algebra}}
\urlstyle{same}
\setlength{\parindent}{0pt}
\setlength{\parskip}{5pt plus 1pt minus 1pt}
\linespread{1.04}
\setlength{\emergencystretch}{3em}
\widowpenalty=10000
\clubpenalty=10000
\displaywidowpenalty=10000
\allowdisplaybreaks[1]
\setlist{nosep,leftmargin=1.5em}
\providecommand{\tightlist}{\setlength{\itemsep}{0pt}\setlength{\parskip}{0pt}}
\providecommand{\pandocbounded}[1]{#1}
\pagestyle{fancy}
\fancyhf{}
\fancyhead[L]{\sffamily\footnotesize\color{muted}OPEN PROBLEMS IN NUMERICAL LINEAR ALGEBRA}
\fancyhead[R]{\sffamily\bfseries\footnotesize\color{accent}IE-20}
\fancyfoot[L]{\parbox[b]{0.9\textwidth}{\sffamily\fontsize{8}{10}\selectfont\color{muted}Editorial ratings; a bounded search cannot certify that a problem remains unsolved.\\Literature check: 2026-09-10}}
\fancyfoot[R]{\sffamily\footnotesize\color{muted}\thepage}
\renewcommand{\headrulewidth}{0.3pt}
\renewcommand{\headrule}{\hbox to\headwidth{\color{accent}\leaders\hrule height \headrulewidth\hfill}}
\makeatletter
\renewcommand{\section}{\Needspace{5\baselineskip}\@startsection{section}{1}{0pt}{12pt plus 2pt minus 2pt}{4pt}{\sffamily\bfseries\large\color{ink}}}
\renewcommand{\subsection}{\Needspace{4\baselineskip}\@startsection{subsection}{2}{0pt}{9pt plus 1pt minus 1pt}{3pt}{\sffamily\bfseries\normalsize\color{ink}}}
\makeatother
\setcounter{secnumdepth}{0}
\begin{document}
{\sffamily\small\color{muted}Linear systems and elimination\par}
\vspace{3pt}
{\raggedright\hyphenpenalty=10000\exhyphenpenalty=10000\sffamily\bfseries\fontsize{21}{24}\selectfont\color{ink}Precision
required for conjugate gradients to attain backward accuracy in n
steps\par}
\vspace{7pt}
{\sffamily\small\textbf{Difficulty:} extreme\quad\textbf{Importance:} interesting
to the community\par}
{\sffamily\small\bfseries\color{accent}Status: Open\par}
\vspace{4pt}
{\color{accent}\hrule height 0.8pt}
\vspace{6pt}
\textbf{Topic:} Finite-precision Krylov methods.

\textbf{Rating rationale:} Extreme reflects a general finite-precision
guarantee for exactly n CG steps, beyond standard exact-arithmetic
termination; community impact is a foundational reliability question for
Krylov solvers.

This fixes an implementation and arithmetic model for the precision
question in the cited workshop report. These conventions are an
editorial specialization of its question, not an additional conjecture
quoted from the authors.

Let \(n\ge1\), \(K\ge1\), and \(0<\varepsilon<1/2\). At precision
\(p\ge2\), the inputs are a real symmetric positive definite matrix
\(A\in\mathbb R^{n\times n}\) and \(b\in\mathbb R^n\setminus\{0\}\)
whose entries are exactly representable using \(p\) binary significant
bits and unbounded exponents. They are read exactly, and the error below
is measured against these stored inputs. Assume \(\kappa_2(A)\le K\).

Use the scalar relative-error model with \(u=2^{-p}\):

\[
\operatorname{fl}(a\circ b)=(a\circ b)(1+\delta),
\quad |\delta|\le u,\quad \circ\in\{+,-,\times,/\},
\]

for nonzero divisors, with no underflow or overflow. Guarantees must
hold for every choice of permitted operation errors. Products are
rounded separately from additions; dot products and matrix-vector rows
are accumulated from left to right in index order, starting from zero.
Copies, sign changes, and comparisons are exact. The model is an error
envelope, not an assumption of independent random roundoff.

Run unpreconditioned conjugate gradients with no restarts, residual
replacement, or reorthogonalization. Initialize \(x_0=0\),
\(r_0=p_0=b\), and \(\rho_0=\operatorname{fl}(r_0^Tr_0)\). For
\(j=0,\ldots,n-1\), evaluate

\[
q_j=\operatorname{fl}(Ap_j),\qquad
d_j=\operatorname{fl}(p_j^Tq_j),\qquad
\alpha_j=\operatorname{fl}(\rho_j/d_j),
\]

\[
x_{j+1}=\operatorname{fl}(x_j+\alpha_jp_j),\qquad
r_{j+1}=\operatorname{fl}(r_j-\alpha_jq_j),
\]

and, if another step is needed,

\[
\rho_{j+1}=\operatorname{fl}(r_{j+1}^Tr_{j+1}),\qquad
\beta_j=\operatorname{fl}(\rho_{j+1}/\rho_j),\qquad
p_{j+1}=\operatorname{fl}(r_{j+1}+\beta_jp_j).
\]

Stop when the computed residual vector is zero, or before a division by
zero. The performance criterion is that at least one iterate produced
before stopping and by step \(n\) satisfies

\[
\eta(x_j;A,b):=
\frac{\|b-Ax_j\|_2}{\|A\|_2\|x_j\|_2+\|b\|_2}\le\varepsilon.
\]

Here \(b-Ax_j\) is the true residual, not the recursively updated vector
\(r_j\). Thus early exact convergence succeeds, while a zero computed
residual or breakdown does not automatically certify success.

\newpage
Define \(P_{\rm CG}(n,K,\varepsilon)\) as the least integer
\(p_{\min}\ge2\) such that this criterion holds at every precision
\(p\ge p_{\min}\), for every admissible stored input at that precision
and every permitted rounding-error sequence; set it to infinity if no
such threshold exists. Determine, up to universal multiplicative
constants, the worst-case dependence of \(P_{\rm CG}\) on \(n\), \(K\),
and \(\varepsilon\).

The question quantifies the precision cost of the exact-arithmetic
finite-termination property for a standard short-recurrence solver.

\subsection{References}\label{references}

\begin{enumerate}
\def\labelenumi{\arabic{enumi}.}
\tightlist
\item
  N. Amsel et al., \href{https://arxiv.org/html/2602.05394v3}{Linear
  Systems and Eigenvalue Problems: Open Questions from a Simons
  Workshop}, arXiv v3 (2026-08-21), §2.6, Problem 2.17. The question
  asks for precision sufficient for normwise backward accuracy in at
  most \(n\) CG steps; neighboring Problems 2.15--2.16 emphasize
  residuals and implementation dependence.
\item
  C. Musco, C. Musco, and A. Sidford,
  \href{https://arxiv.org/abs/1708.07788}{Stability of the Lanczos
  Method for Matrix Function Approximation}, SODA 2018, 1605--1624,
  abstract and discussion of inverse approximation and finite-precision
  CG. Its polynomial-approximation results are relevant background, not
  a sharp answer to this target.
\item
  S. Chenakkod, M. Dereziński, X. Dong, and M. Rudelson,
  \href{https://arxiv.org/abs/2604.23193}{Well-Conditioned Oblivious
  Perturbations in Linear Space}, 2026, abstract and perturbed-CG
  application.
\end{enumerate}

\subsection{Status check --- 2026-09-08}\label{status-check-2026-09-08}

Checked workshop v3, which explicitly retains Problem 2.17 in its August
2026 update. Searches combined ``conjugate gradient'', ``precision'',
``bits'', ``n steps'', ``backward error'', and 2026. No sharp solution
was located. The April 2026 perturbation paper concerns a modified
randomized solver and an \(O(n)\) matrix-vector bound, not this fixed
unperturbed recurrence with an at-most-\(n\) requirement. Input rounding
is excluded explicitly; other CG implementations and stronger bit-cost
claims are outside this model specialization. The bounded literature
search is not a proof of current openness.

\subsection{Audit update --- 2026-09-10}\label{audit-update-2026-09-10}

Rechecked \href{https://arxiv.org/html/2602.05394v3}{the August 2026
workshop revision}, Problem 2.17. It still asks the required precision
question. CG backward-accuracy and finite-precision searches located no
sharp answer for the explicitly fixed arithmetic model here; those
implementation conventions remain an editorial specialization rather
than a verbatim source claim.
\end{document}
